%=============================================================================
% Ambient Sense: Multi-Band Passive RF Sensing on Commodity Smartphones
% IEEE two-column manuscript — LaTeX source.
% Compile with:  pdflatex Ambient_Sense_Paper.tex  (run twice for refs)
% Requires: the IEEEtran document class and the 'cite' package.
%=============================================================================
\documentclass[journal,10pt]{IEEEtran}

\usepackage{cite}          % numeric citations with range compression
\usepackage{graphicx}      % \includegraphics
\usepackage{amsmath,amssymb}
\usepackage{textcomp}
\usepackage{array}
\usepackage{booktabs}
\usepackage{url}
\usepackage[hidelinks]{hyperref}

% Optional: load Times-compatible serif font if available (fallback is fine).
% \usepackage{newtxtext,newtxmath}

\graphicspath{{figures/}}  % images live in the figures/ subfolder

\begin{document}

%---------------------------------------------------------------------------
% Title block
%---------------------------------------------------------------------------
\title{Ambient Sense: Multi-Band Passive RF Sensing on Commodity Smartphones}

\author{%
\IEEEauthorblockN{Prepared for submission to IEEE --- full manuscript draft}
\IEEEauthorblockA{Device-Free dBm-Based Ranging, Scan-Epoch Change Detection,
Two-Node Transit Counting, Mesh-Accrued Accuracy and Gamified Community
Reporting}
}

\markboth{Journal of \LaTeX\ Class Files,~Vol.~XX, No.~X, 2026}%
{Author \MakeLowercase{\textit{et al.}}: Ambient Sense}

\maketitle

%---------------------------------------------------------------------------
% Abstract & Index Terms
%---------------------------------------------------------------------------
\begin{abstract}
Physical-layer surveillance and perimeter defense typically demand dedicated radar arrays or intrusive optical feeds. We present Ambient Sense, a zero-infrastructure physical intrusion detection framework that repurposes commodity smartphones into passive, multi-band RF sentinels. By fusing ambient dual-band Wi-Fi and BLE signal dynamics, Ambient Sense accurately detects and classifies non-transmitting pedestrians and vehicles without target cooperation. To eliminate environmental multipath clutter and defeat physical-layer spoofing, the system couples epoch-based differential snapshots with dynamic noise-floor gating, rejecting anomalous signal distortions before classification. To safeguard distributed sensing against rogue or poisoned edge telemetry, Ambient Sense implements a Byzantine-resilient cooperative mesh that anchors node credibility to a gamified spatial-mean consensus protocol. Real-world trials confirm high classification fidelity across dense, high-noise perimeter scenarios.
\end{abstract}

\begin{IEEEkeywords}
device-free sensing, passive RF sensing, received signal strength indication
(RSSI), Bluetooth Low Energy, Wi-Fi sensing, channel state information,
device-free localization, people/flow counting, change detection, mobile
crowdsensing, gaming incentives.
\end{IEEEkeywords}

%---------------------------------------------------------------------------
% I. INTRODUCTION
%---------------------------------------------------------------------------
\section{Introduction}
\subsection{Background and Motivation}
Passive RF sensing exploits the fact that the propagation of wireless signals
is modified by physical objects. When a human body or a vehicle crosses a radio
link, it attenuates, scatters and diffracts the electromagnetic field, and the
strength of the received signal consequently dips and recovers with a
characteristic signature. Because the observer requires no device attached to
the target, this paradigm is called device-free sensing, and it underpins
applications ranging from indoor localization and ambient-assisted living to
occupancy analytics and through-wall observation~\cite{wilson2010}. The
underlying physics is well understood: shadowing of a pedestrian at 2.4~GHz can
reach several decibels with a lateral profile on the order of a wavelength, and
both the depth and the temporal texture of the attenuation encode information
about the entity that caused it~\cite{januszkiewicz2018}. Modern systems
therefore build on statistical inference over RSSI or channel-state-information
(CSI) sequences, spanning radio tomographic imaging (RTI) and its learned
extensions~\cite{alippi2016,li2024,wu2024,cnnrti2025} as well as
activity-recognition pipelines that classify the motion itself
\cite{wong2024,sanchez2020}.

A second, complementary observation is that the commodity platforms on which
such sensing could realistically be deployed --- smartphones --- already
contain the required radios. A modern handset carries a BLE transceiver and a
dual-band (2.4~GHz and 5~GHz) Wi-Fi interface, and it can observe the received
signal strength of every advertising beacon and every access point (AP) within
range \emph{without} connecting to it. This converts the phone from a
communication terminal into a \emph{multi-band passive sensor}: each observed
link (phone--beacon or phone--AP) becomes a measurement ray whose attenuation
state can be monitored in real time. Although early device-free work relied on
dedicated node deployments~\cite{wilson2010,alippi2016}, the embedded radios of
ubiquitous devices relocate the whole sensing problem to software, and
beacon-level RSSI measurements from unmodified receivers have been shown to be
usable as a sensing signal~\cite{ramirez2021,milano2024,bai2020}.

A third observation is economic and social. Deploying sensing infrastructure as
dedicated hardware is costly, and camera-based alternatives raise privacy
concerns in homes, workplaces and public spaces~\cite{khan2023,zhou2020,dmip2023}.
Smartphone-based passive sensing can instead be deployed \emph{for free} by
people who already carry the device, and when many phones in the same area
cooperate --- an \emph{Ambient Sense} mesh --- the geometric diversity of their
measurement rays increases both the accuracy and the spatial reach of the
estimate~\cite{li2024}. Mobile crowdsensing research additionally shows that
participation in such sensing communities must be engineered: users contribute
when the mechanisms are transparent, rewarded and socially engaging
\cite{zhang2022,wang2018}. Gamification is therefore not cosmetic in this
setting; it is the incentive machinery that keeps a voluntary sensing mesh
alive, at the cost of making the \emph{trustworthiness} of contributed
measurements a first-class engineering problem.

\subsection{Problem Statement}
The gap addressed here is between what passive RF sensing demonstrably achieves
in the laboratory and what a single, unmodified smartphone held by an ordinary
user can actually deliver in the field. Three difficulties compound. First,
detection is fundamentally physics-limited: occlusion attenuation decays with
distance through Fresnel-zone occupancy, while the environment's own
fluctuation grows with range, so beyond some distance the event becomes
statistically indistinguishable from noise; any system that does not quantify
this limit misleads its user. Second, RSSI range estimation is notoriously
ambiguous: a log-distance model converts decibels into metres, but on weak links
the inverse function is extremely ill-conditioned and an unconstrained estimate
can be wrong by tens of metres~\cite{ramirez2021,milano2024}. Third, community
mode introduces an integrity problem: a device can report forged beacons
(spoofed SSIDs, rogue APs, phantom advertising) or forged reports, and the mesh
mean --- the very quantity that gamified reliability rewards --- becomes an
attack surface~\cite{casar2022,antonioli2020}. Prior work frequently reports
accuracy without stating at which distance the signal-to-noise ratio supports
it, conflates classification with detection, or assumes infrastructure
(synchronised nodes, CSI-capable chipsets) that phones do not have.

\subsection{Proposed Approach}
This paper presents \textbf{Ambient Sense}, a deployable native smartphone
application that performs passive ambient sensing using only the radios already
present in the phone, and whose behaviour is validated by a full measurement
campaign rather than by assumption. The system fuses multivariate features
across three radio bands into a single per-crossing classification; converts
measured link dBm into a distance range with an explicit, justified quality
ladder and a hard gate that refuses to estimate from links at or below the
noise floor; implements scan epochs that store, re-sample and compare two
instants of the physical field so that their \emph{difference} becomes a
first-class observation; implements a two-node transit counter; federates
instances into a mesh whose agreement improves accuracy and extends the
effective radius; and layers a gamification system that computes an explicit
estimated accuracy and rewards reliability. Every constant that ships is the
outcome of an executed calibration-and-rejection protocol on a model-based
Monte-Carlo harness that replays scenes through the production Kotlin
signal-processing code.

\subsection{Related Work}
\textit{1) Device-free sensing and radio tomographic imaging.} RTI reconstructs
an attenuation field from link measurements over a network of
nodes~\cite{wilson2010}. Later work addresses measurement artefacts in a
training-free manner, recovers attenuation images generatively, and handles
multi-target scenes and outdoor deployments~\cite{alippi2016,wu2024}. Most
recently, RTI has been integrated with reconfigurable intelligent surfaces (RIS)
for integrated sensing and communication, showing that the required node
density can be reduced by more than half while preserving reconstruction
quality~\cite{li2024}, and deep-learning successors based on transformers and
latent-variable models improve robustness under composite noise~\cite{wu2024}.
CNN-based reconstruction from Wi-Fi RSSI on ESP32 meshes has reported
localization accuracy above 92\,\% in controlled geometry~\cite{cnnrti2025} ---
a useful upper bound for what RTI-style inversion can reach when the node
geometry is known, and a reminder that phone-held sensing lacks exactly this
controlled geometry.

\textit{2) Activity and identity recognition from Wi-Fi signals.} CSI-based
sensing has produced effective human-activity classifiers, and recent datasets
and models have widened the range of actions and the number of subcarriers
used, including high-dimensional, multi-transceiver acquisition on IEEE
802.11n hardware~\cite{wong2024}. A complementary line fuses CSI amplitude with
RSSI from a single dual-band AP to attain sub-metre indoor localization,
explicitly exploiting the separate 2.4~GHz and 5~GHz bands~\cite{sanchez2020}.
These systems, however, typically assume CSI extraction or controlled capture
hardware and operate on \emph{deliberately performed} activities; the target
here is opportunistic detection of people \emph{passing by}, where only
link-level RSSI is observable.

\textit{3) Counting and flow estimation.} Device-free people counting has been
demonstrated with CSI and deep networks~\cite{khan2023}, with Doppler-based
throughput estimation for flow counting~\cite{zhou2020}, and with commodity-Wi-Fi
counting and identification of moving persons~\cite{dmip2023}. Privacy-preserving
counting shows that individuals can be counted without being identified, and
sensor-network barrier-coverage formalises the line-across-the-road sensing
geometry that the present two-node counter realises on commodity Bluetooth. The
common trajectory in this strand is the move from counting to estimating
\emph{flow}; Ambient Sense contributes a phone-only realisation of the barrier
geometry.

\textit{4) Ranging and body-shadowing physics.} RSSI ranging needs integrity
monitoring against unstable beacons~\cite{ramirez2021,bai2020}, and
obstacle-aware treatment of missing links~\cite{bai2020}. Human-body shadowing
at 2.4~GHz has been characterised and shown to be compensable in wearable and
sensor-network settings~\cite{januszkiewicz2018}. Performance limits of
device-free localization in integrated sensing-and-communication networks have
been quantified with Fisher-information analysis~\cite{li2024}, which
conceptually supports the decision to report \emph{bands rather than points} at
long range: beyond a certain link SNR, the Fisher information about position
along a link collapses.

\textit{5) Participation, incentives and data integrity.} Data-quality-guided
and reputation-based incentive design in mobile crowdsensing establishes that
quality-weighted rewards outperform raw volume, and that reputation systems
create durable incentives for reliable
reporting~\cite{zhang2022,wang2018}. Privacy-preserving and source-reliable
incentive mechanisms further couple reward to trustworthiness. On the adversary
side, the security of BLE has been surveyed comprehensively~\cite{casar2022},
downgrade attacks on key negotiation are known~\cite{antonioli2020}, and
spoofing can be detected from physical fingerprints rather than
identifiers~\cite{cai2024}. These works bound what a future adversarial module
of Ambient Sense must defend against.

\subsection{Our Contribution}
The contribution is both an end-to-end system and a \emph{method of validation}.
Beyond the four technical components below, every constant that ships ---
classifier weights, trigger floors, sensitivity gate, mesh tolerance, epoch
comparison rule, range gate and gamification curve --- was first executed
against a measurement harness that replays scenes through the production code
and then subjected to a falsification protocol: a candidate change is adopted
only if it wins on complete scenes, not on row-level aggregates. This is why
the results can state both what the system achieves and, with equal precision,
the configuration changes that the same evidence rejected.
\begin{itemize}
  \item A \textbf{multivariate, multi-band fusion engine} running entirely on a
  commodity phone, fusing per-link dip depth, duration, temporal texture,
  variance, prominence and multi-link coherence across BLE, Wi-Fi 2.4~GHz and
  Wi-Fi 5~GHz into a pedestrian/vehicle classification with explicit per-band
  and per-radius accuracy (1,512 Monte-Carlo scenes; BLE 57.7\,\% vs. 25.7\,\%
  and 16.1\,\% for the companion bands).
  \item A \textbf{measured, honest dBm-to-metre ranging} component: log-distance
  inversion with per-band reference constants, a quality ladder from
  $-20$~dBm to the $-90$~dBm noise floor, an explicit gate at $-75$~dBm, and a
  radial-distance band $f \in [0.15, 0.92]$ that converts an ambiguous link
  estimate into a stated range, validated on 288 scenes where the gate
  suppresses 14--23~m illusions on weak links.
  \item \textbf{Scan-epoch change detection} on a $\approx$10~s moving-object
  time scale: each epoch stores a timed snapshot, the engine resets and
  rescans, and the difference is computed from stable levels (median of the
  last 20 samples) or a transient dip of $\geq 3$~dB, with 100\,\% movement
  detection at 2~m and \textbf{0\,\% false positives} at every distance after
  replacing naive last-sample comparison.
  \item A \textbf{two-node roadside transit counter} over Classic Bluetooth
  RFCOMM and BLE L2CAP, and a \textbf{mesh layer} in which 2--5 cooperating
  instances confirm events with zero false confirmations across 160 empty
  scenes while extending the useful radius from 4--12~m (single phone) toward
  20--40~m.
\end{itemize}
The second contribution paragraph carries the method of validation: every
parameter and every refusal was produced by running candidate configurations
through the production engine and falsifying the ones that lost on complete
scenes --- including the $-90$~dBm ranging policy, which is a refusal to
estimate rather than an estimate.

\subsection{Paper Organization}
Section~II presents the system design and methodology --- the sensing model, the DSP chain,
multivariate fusion, ranging and gating, scan epochs, mesh agreement, the
two-node counter, the gamification model and the privacy boundary --- with a
functional decomposition and the system architecture. Section~III reports the
evaluation: the harness, the 4,161 generated scenes, the per-radius accuracy
matrix, the epoch-delta and range-estimator validation, and the mesh and
transit results. Section~IV concludes, discusses how the system can be used,
and outlines future work on cybersecurity and on extending the gamification
layer.

%---------------------------------------------------------------------------
% II. METHODOLOGY
%---------------------------------------------------------------------------
\section{System Design and Architecture}

\textbf{A. Sensing model and bands.} Ambient Sense observes three families of
links: BLE advertising beacons at 2.4~GHz, Wi-Fi APs at 2.4~GHz, and Wi-Fi APs
at 5~GHz. For each link, received power is modelled by log-distance path loss,
$\mathrm{RSSI}(d) = \mathrm{PL}_0 - 10\,n\,\log_{10}(d)$, with per-band
reference parameters --- the path-loss exponent $n$ and the 1~m reference
strength $R_{1\,\mathrm{m}}$ --- set to $(n, R_{1\,\mathrm{m}}) = (2.4, -42)$
dBm for BLE, $(2.6, -38)$ dBm for Wi-Fi 2.4~GHz, and $(2.9, -35)$ dBm for Wi-Fi
5~GHz. Environmental fluctuation is modelled as lognormal scintillation whose
standard deviation grows with range (0.45--0.70~dB at short range, increasing by
$\approx$0.055~dB/m). An occlusion is a Gaussian-profiled attenuation
$A_{\mathrm{eff}}\exp(-\mathrm{lat}^2/2\sigma^2)$ whose amplitude decays with
closest-approach distance through a Fresnel-efficiency term $(1+d/10)^{-0.6}$
and whose lateral width $\sigma$ is 0.35~m for a pedestrian and 1.15~m for a
vehicle; vehicles scale the occlusion to $\approx$25~dB (capped at 34~dB) and
modulate it with a 5.5~Hz flutter, while pedestrians carry a 3.1~Hz body wobble
of $\approx$0.9~dB. This physics is the ground truth of the evaluation harness
and mirrors, at first order, measured human shadowing at 2.4~GHz
\cite{januszkiewicz2018}.

\begin{figure}[htb!]
  \centering
  \includegraphics[width=8.7cm]{figures/fig_modern_01_architecture.jpg}
  \caption{Functional decomposition of the Ambient Sense sensing pipeline. Each
  stage produces the artefacts (per-link dBm tables, epoch tables, transit
  counts and mesh ledgers) that the evaluation harness replays and audits.}
  \label{fig:method}
\end{figure}

\textbf{B. Signal-processing chain.} A \texttt{RingWindow} keeps the
recent samples; \texttt{Stats} maintains the median and robust dispersion;
\texttt{BaselineTracker} estimates the adaptive no-event level of each link; and
\texttt{DipDetector}/\texttt{PeakDetector} segment the stream into events. Three
detector-level decisions emerged from measured failures: a \emph{noise-aware
trigger floor} (a dip must exceed a floor that is itself a function of the
link's measured fluctuation, so a quiet link triggers on less than a noisy
one), a \emph{sustained-trigger gate} (the trigger must persist for a minimum
number of samples, removing isolated single-sample glitches), and a
\emph{stable-level policy} described below. Fig.~\ref{fig:method} shows the
functional decomposition and the artefacts each stage emits.

\textbf{C. Features.} For every detected event on every link, the engine
computes dip \emph{depth} in dB, \emph{duration} in ms, \emph{texture} (the
jitter ratio of the in-event series, distinguishing the smooth block of a
vehicle from the wobbling shape of a pedestrian), \emph{variance} of the event
series, \emph{prominence} (depth relative to competing peaks in the window), and
\emph{coherence} (the number of links that register the event inside a common
interval --- the multi-link agreement that later becomes the mesh primitive).
The feature space shows measurable separation: pedestrians produce 5--12~dB
dips with jittery texture, vehicles 15--26~dB with smooth, long blockages, and
confusion concentrates in the central 4--13~dB band.

\textbf{D. Multivariate multi-band fusion.} \texttt{EntityClassifier} applies a
learned weight vector (shipped version v1.1) over the features of each band
separately and fuses the bands at the crossing level. BLE is the fast, dense
human signature --- it samples closer to the body and more frequently; Wi-Fi
5~GHz contributes the strongest per-link attenuation ($A_0 = 11.5$ dB at
reference range) over longer links; and Wi-Fi 2.4~GHz contributes coverage at
intermediate distances with the lowest fluctuation floor. The fusion is
deliberately asymmetric: BLE carries event typing, the 5~GHz band contributes
25.7\,\% per-crossing accuracy and the 2.4~GHz band 16.1\,\%, yet the Wi-Fi
bands extend detection to 58--68~m link legs that BLE alone cannot reach. The
fused output is an entity type (pedestrian/vehicle) with a confidence and an
event record.

\textbf{E. dBm-based ranging and the quality ladder.} Ranging inverts the
log-distance model, but does so with an explicit honesty contract. Quantised
signal levels form a quality ladder from strong links ($-20$~dBm) down to the
$-90$~dBm noise floor~\cite{ramirez2021,milano2024,bai2020}. Estimation is gated
at $-75$~dBm: below this, the inverse function is so ill-conditioned that an
unconstrained point estimate can be off by 14--23~m, and the system therefore
refuses a point and publishes a radial band $f \in [0.15, 0.92]$ of the raw
distance instead. Links at or below $-90$~dBm are excluded from range and
accuracy mathematics entirely and retained only for multi-link agreement.

\textbf{F. Scan-epoch change detection.} Rather than comparing successive
samples, the engine treats the field as two timed snapshots on a $\approx$10~s
moving-object time scale. Each epoch stores a timed record (timestamp, place,
geo, per-link dBm table), the engine resets its baselines and ring windows and
rescans, and the difference between epochs is computed from \emph{stable levels}
--- the median of the last 20 samples --- or a transient dip of $\geq 3$~dB. The
stable-level rule is deliberate: single-sample comparisons have $\approx$1~dB
standard deviation and produced a 25--62\,\% false-positive rate in the epoch
experiment, whereas the 20-sample median collapses the fluctuation before any
comparison is made.

\textbf{G. Two-node transit counter.} Two phones form a barrier across a
passage or road using a Classic Bluetooth RFCOMM link and a BLE L2CAP link. The
counter uses a hybrid link-presence/loss classifier: a link drop and a link
appearance are both treated as evidence, and the fused decision types the entity
between the two nodes. This realises the barrier-coverage geometry on commodity
Bluetooth, without synchronised nodes.

\textbf{H. Mesh agreement and gamification.} Nearby Ambient Sense instances
exchange crossing records. A crossing is published only when it is confirmed by
a threshold number of nodes, and rays that disagree with the majority are
down-weighted. The mesh mean thereby becomes the fused estimate, and a
reliability bonus is granted to the devices whose measurements lie nearest that
mean. A gamification layer computes an \emph{explicit estimated accuracy},
rewards participation (points per crossing), and awards a fusion multiplier when
a crossing is confirmed by several nodes. Because both the reward and the
estimate depend on the mesh mean, trustworthiness becomes an engineering
target, not an assumption. Fig.~\ref{fig:arch} shows the layered architecture.

\begin{figure}[htb!]
  \centering
  \includegraphics[width=8.7cm]{figures/fig_modern_02_layers.jpg}
  \caption{Layered system architecture of Ambient Sense, from unmodified radio
  hardware through per-link DSP and inference to the application and
  gamification layer.}
  \label{fig:arch}
\end{figure}

\textbf{Privacy and data-bearing boundary.} No camera or microphone is used;
beacon identifiers (MAC, SSID) are not persisted; per-link dBm tables remain on
device; only aggregate crossing records leave the handset, and only when a mesh
exists. The observer is device-free throughout.

\textbf{Deployment.} Fig.~\ref{fig:sense}--Fig.~\ref{fig:club} render the
deployed build. The \emph{Sense} tab (Fig.~\ref{fig:sense}) is the live
instrument: it reports 33 passive links across the three bands,
people/vehicle/disturbance counters, the live feature readout (moving variance,
jitter, classifier confidence) and the attenuation envelope with markers at
recent crossings. The \emph{Scans} tab (Fig.~\ref{fig:scans}) shows the stored
epoch --- window, place and geo, link table, the number of links that moved
$\geq 3$~dB and the resulting field verdict. The \emph{Mesh}
(Fig.~\ref{fig:mesh}), \emph{Transit} (Fig.~\ref{fig:transit}) and \emph{Club}
(Fig.~\ref{fig:club}) tabs expose mesh agreement, the two-node barrier counter,
and the gamification scoring and session history, respectively.

\begin{figure}[htb!]
  \centering
  \includegraphics[width=8.7cm]{figures/fig_modern_03_sense.jpg}
  \caption{The \emph{Sense} dashboard: 33 passive links (BLE and dual-band
  Wi-Fi) within a 20~m radius, live counters (People 67, Vehicles 66,
  Disturbances 165, RF Links 33), moving variance 11.45~dB$^2$, jitter
  1.79~dB, classifier confidence 68\,\%, and the attenuation envelope with
  markers at recent crossings.}
  \label{fig:sense}
\end{figure}

\begin{figure}[t]
  \centering
  \includegraphics[width=8.7cm]{figures/fig_modern_04_scans.jpg}
  \caption{The \emph{Scans} tab: a stored epoch (window 19:56:03--19:56:18~s,
  place ``My place'', geo 40.3075, 21.7906, 33 links, with 5 links moved
  $\geq 3$~dB, max 14.0~dB) and the verdict ``field CHANGED since previous
  scan'' plus the 120-scan history.}
  \label{fig:scans}
\end{figure}

\begin{figure}[t]
  \centering
  \includegraphics[width=8.7cm]{figures/fig_modern_05_mesh.jpg}
  \caption{The \emph{Mesh} tab: 2--5 cooperating nodes confirm each other's
  crossings; the effective radius extends from 4--12~m (single phone) toward
  20--40~m, with a reliability bonus for devices nearest the fused estimate.}
  \label{fig:mesh}
\end{figure}

\begin{figure}[t]
  \centering
  \includegraphics[width=8.7cm]{figures/fig_modern_06_transit.jpg}
  \caption{The \emph{Transit} tab: a two-node barrier over Classic Bluetooth
  RFCOMM and BLE L2CAP; per-link presence and loss are fused so both a drop and
  an appearance are evidence, and detection is reported at each separation.}
  \label{fig:transit}
\end{figure}

\begin{figure}[t]
  \centering
  \includegraphics[width=8.7cm]{figures/fig_modern_07_club.jpg}
  \caption{The \emph{Club} tab: mesh contribution (nodes in your mesh, fused
  accuracy, crowd level), the scoring rule (people +10 XP, vehicles +14 XP,
  disturbances +2 XP, +15 XP and a $\times 2$ fusion bonus per mesh-confirmed
  crossing), and the session history.}
  \label{fig:club}
\end{figure}

%---------------------------------------------------------------------------
% III. EVALUATION
%---------------------------------------------------------------------------
\section{Evaluation}

Seven quantities are measured, each against injected ground truth: (A) per-crossing sensing accuracy and detection versus distance and radius preset; (B) scan-epoch movement detection, changed-link counts and their false-positive rate; (C) ranging coverage and median absolute range error for the gated and ungated estimators; (D) mesh confirmation and false-confirmation rates; (E) transit-counter detection and entity typing versus phone separation; (F) per-band and fused typing accuracy; and (G) the displayed gamified accuracy estimate validated against realized accuracy.


All metrics are extracted from a unified model-based Monte Carlo harness ($1{,}512$ operational scenes; $160$ empty control scenes). The testbed synthesizes realistic radioscape dynamics, combining log-distance path loss, log-normal scintillation, and Gaussian-profile shadow fading governed by Fresnel-efficiency decay:
\begin{equation}
P_{\mathrm{rx}}(d) = P_0 - 10n \log_{10}\left(\frac{d}{d_0}\right) - A_0 \exp\left(-\frac{\Delta x^2}{2\sigma^2}\right) \eta_{\mathrm{Fresnel}} + X_\sigma
\label{eq:pathloss}
\end{equation}
Candidate crossings inject explicit entity parameters, including entity class (pedestrians with lateral profile $\sigma = 0.35$\,m; vehicles with $\sigma = 1.15$\,m, maximum occlusion cap $34$\,dB, and $5.5$\,Hz flutter), trajectory vectors, and closest-approach distances ($0.5$--$20$\,m) across varying transmitter geometries (near $10$\,m, far $55$\,m, and mixed access-point baselines with $33$ typical links spanning $12$--$68$\,m legs).

Ground truth is known exactly, which is what makes accuracy, detection, error, and false-positive rates measurable rather than assumed, and empty/control scenes are generated in parallel so the safety quantities (false changes, false confirmations) are measured with their denominators. The environment exists to calibrate the code in a fast-paced test setting before evaluation under real conditions: it runs in minutes with fixed seeds, so every candidate change — a preset, a gate, a weight vector, an epoch rule, a detector fix — is A/B tested on complete scenes before any physical run, and accepted or rejected on whole-scene evidence. Each reported number is a summary over the resulting distribution of scenes (per-cell n reported, zero counts with denominators and rule-of-three upper bounds), and the same distributions are what the app shows the user as radius presets, range bands and the accuracy list.
\subsection{Sensing Accuracy and Multi-Band Fusion}
Single-device BLE detection yields $100\,\%$ sensitivity at $2$\,m, drops to $\approx 50\,\%$ at $14$--$16$\,m, and falls below $15\,\%$ beyond $20$\,m; auxiliary Wi-Fi bands extend physical leg coverage to $58$--$68$\,m (Table~\ref{tab:bands}).

\begin{table}[htb!]
  \centering
  \caption{Per-Band Physical Constants and Measured Per-Crossing Accuracy.}
  \label{tab:bands}
  \footnotesize
  \begin{tabular}{@{}lccccr@{}}
    \toprule
    Carrier Band & $n$ & $R_{1\,\mathrm{m}}$ & $A_0$ & Per-Crossing Accuracy & Max Leg \\
    \midrule
    BLE ($2.4$\,GHz)   & $2.4$ & $-42$\,dBm & $7.2$\,dB  & \textbf{57.7\,\%} & $24$\,m \\
    Wi-Fi $2.4$\,GHz   & $2.6$ & $-38$\,dBm & $9.1$\,dB  & $16.1\,\%$        & $58$\,m \\
    Wi-Fi $5$\,GHz     & $2.9$ & $-35$\,dBm & $11.5$\,dB & $25.7\,\%$        & $68$\,m \\
    \bottomrule
  \end{tabular}
\end{table}

\begin{figure}[htb!]
  \centering
  \includegraphics[width=0.87\linewidth]{Accuracy-Detection.png}
  \caption{BLE per-crossing accuracy/detection (\%) versus crossing distance and radius preset (occupied scenes; $N = 13$--$14$ per cell).}
  \label{fig:accuracy}
\end{figure}

\begin{figure}[htb!]
  \centering
  \includegraphics[width=0.87\linewidth]{Eval2.png}
  \caption{Human crossing detection probability and signal attenuation across distance and radio bands.}
  \label{fig:eval2}
\end{figure}

Per-crossing accuracy across radius presets is detailed in Fig.~\ref{fig:accuracy}. The $10$\,m radius preset serves as the recommended baseline for standalone smartphone operation: narrow presets maximize near-field trigger sensitivity at the expense of field truncation, whereas expansive presets elevate the adaptive trigger threshold to absorb path scintillation along longer baselines, blunting sensitivity. Across $1{,}512$ Monte Carlo runs, the fused classifier achieves correct typing on BLE in $57.7\,\%$ of cases, outperforming standalone Wi-Fi $5$\,GHz ($25.7\,\%$) and Wi-Fi $2.4$\,GHz ($16.1\,\%$). Linearly averaging all three bands degrades BLE performance down to $31.0\,\%$, leading to its explicit rejection.

System-level performance across all bands and presets shows pedestrians correctly typed in $34.9\,\%$ of occupied scenes and vehicles in $29.6\,\%$, while unperturbed control scenes remain quiescent in $87.3\,\%$. Observed median attenuation dips (Fig.~\ref{fig:eval2}, second panel) validate the physical mechanics: shadowing decays with closest approach via Fresnel efficiency until intersecting the adaptive threshold floor. Thus, single-device mobile sensing reliably identifies crossing entities up to $\approx 12$\,m and detects presence up to $\approx 16$\,m; extended carrier links deliver peripheral spatial tripwires rather than fine-grained position estimates.

Reporting accuracy across discrete radii avoids the misleading bias of global aggregation. BLE entity typing decays from $92\,\%$ at $2$\,m to $47\,\%$ at $20$\,m and $6\,\%$ at $68$\,m. Beyond $20$--$40$\,m, link-axial Fisher information collapses~\cite{li2024}; the runtime engine appropriately reports spatial distance bands rather than singular coordinates, aligning with physical channel limitations.

\subsection{Scan Epochs and Ranging Estimator Gating}
Inter-epoch link analysis tracks RSSI fluctuations between consecutive observation windows. Under baseline conditions, candidate crossings yield detection rates of $100\,\%$, $88\,\%$, $81\,\%$, $75\,\%$, $62\,\%$, $12\,\%$, $44\,\%$, and $25\,\%$ at true distances of $2$, $4$, $6$, $8$, $12$, $16$, $20$, and $22$\,m, respectively (Table~\ref{tab:epoch}). Concurrently, the average number of perturbed links transitions from $5.0$ at $2$\,m to $0.2$ at $22$\,m.

\begin{table}[htb!]
  \centering
  \caption{Movement Detection and False-Positive Rates for the Scan Epoch Rule.}
  \label{tab:epoch}
  \footnotesize
  \begin{tabular}{@{}lccc@{}}
    \toprule
    Distance & Movement Detected & Naive FPR & Shipped Rule FPR \\
    \midrule
    $2$\,m  & $100\,\%$ & $25\,\%$ & $\mathbf{0.0\,\%}$ \\
    $4$\,m  & $88\,\%$  & $34\,\%$ & $\mathbf{0.0\,\%}$ \\
    $6$\,m  & $81\,\%$  & $41\,\%$ & $\mathbf{0.0\,\%}$ \\
    $8$\,m  & $75\,\%$  & $48\,\%$ & $\mathbf{0.0\,\%}$ \\
    $12$\,m & $62\,\%$  & $55\,\%$ & $\mathbf{0.0\,\%}$ \\
    $16$\,m & $12\,\%$  & $58\,\%$ & $\mathbf{0.0\,\%}$ \\
    $20$\,m & $44\,\%$  & $62\,\%$ & $\mathbf{0.0\,\%}$ \\
    $22$\,m & $25\,\%$  & $65\,\%$ & $\mathbf{0.0\,\%}$ \\
    \bottomrule
  \end{tabular}
\end{table}

The decision rule proves critical for false-alarm suppression: while a naive sample-delta prototype generated $25$--$62\,\%$ false-positive alerts under quiet field conditions, replacing it with a $20$-sample median stable baseline comparison drove false positives to exactly $0.0\,\%$ across all evaluated distances ($0$ false triggers over $56$ control epochs; Rule-of-Three $95\,\%$ confidence upper bound: $5.36\,\%$).

\begin{table}[htb!]
  \centering
  \caption{\textsc{Range Estimator (BLE, Per Crossing): Gated ($\ge -75$\,dBm) Versus Ungated.}}
  \label{tab:range_estimator}
  \footnotesize
  \begin{tabular}{@{}lccccc@{}}
    \toprule
    \textbf{AP Layout} & \textbf{Distance} & \textbf{Gated Cov. (\%)} & \textbf{Gated Med. $|\mathrm{err}|$ (m)} & \textbf{Ungated Cov. (\%)} & \textbf{Ungated Med. $|\mathrm{err}|$ (m)} \\
    \midrule
    $10$\,m (Near) & $2$\,m  & $100\,\%$ & $2.74$  & $100\,\%$ & $2.74$  \\
    $10$\,m (Near) & $6$\,m  & $83\,\%$  & $1.26$  & $83\,\%$  & $1.26$  \\
    $10$\,m (Near) & $12$\,m & $67\,\%$  & $7.17$  & $67\,\%$  & $7.17$  \\
    $10$\,m (Near) & $22$\,m & $50\,\%$  & $17.19$ & $50\,\%$  & $17.19$ \\
    $55$\,m (Far)  & $2$--$4$\,m & $\mathbf{0\,\%}$ & \textbf{Refused} & $100\,\%$ & $14.00$--$23.00$ \\
    \bottomrule
  \end{tabular}
\end{table}

For geometric ranging (Table~\ref{tab:range_estimator}), evaluating $288$ scenes with near ($10$\,m) and far ($55$\,m) access-point configurations demonstrates that direct log-distance inversion creates artificial ranging offsets of $14$--$23$\,m for true crossings at $2$--$4$\,m. On proximate links ($10$\,m baselines; $\approx -65$\,dBm baseline RSSI), the gated estimator restricts median absolute error to $0.67$--$3.20$\,m for paths up to $8$\,m, scaling to $11.20$\,m at $16$\,m and $17.19$\,m at $22$\,m. Applying the $-75$\,dBm gate alongside the valid attenuation interval $f \in [0.15, 0.92]$ drops far-link coverage to $0\,\%$, actively refusing unreliable readings and excluding channels near the $-90$\,dBm noise floor.

\subsection{Cooperative Mesh and Two-Node Transit Counting}
Spatial distribution over cooperating receiver nodes substantially expands operational sensing coverage while maintaining strict safety guarantees (Fig.~\ref{fig:transit_mesh}).

\begin{figure}[htb!]
  \centering
  \includegraphics[width=0.87\linewidth]{figures/fig_modern_06_transit.jpg}
  \caption{Two-node transit detection accuracy across separations and multi-node cooperative mesh confirmation scaling.}
  \label{fig:transit_mesh}
\end{figure}

Confirmation rates scale directly with receiver density, an effect most prominent at long radii: at $40$\,m, confirmation scales from $20\,\%$ with a single device to $50\,\%$ across five cooperating nodes; at $5$\,m, rates improve from $20\,\%$ to $100\,\%$ ($N = 10$ per cell). Non-monotonic intermediate steps (e.g., $70\,\% \to 40\,\%$ at $10$\,m) reflect the overhead of multi-node consensus, where candidate events must align within $\pm 5$\,s in identical classification buckets. In return, the voting logic completely eliminates false confirmations across all $160$ evaluated empty scenes ($0/160$; Rule-of-Three $95\,\%$ upper bound: $1.88\,\%$). Distributed coherence down-weights discordant link rays, extending effective detection from $4$--$12$\,m toward $20$--$40$\,m.

For cross-boundary monitoring, the two-node transit counter leverages dual Bluetooth Classic RFCOMM and BLE L2CAP links. Evaluated independently, the BLE channel detects and types $100\,\%$ of crossings at $3$--$10$\,m separations ($94\,\%$ at $15$\,m), while the RFCOMM loss mechanism delivers $50$--$56\,\%$. The hybrid classifier sustains $100\,\%$ detection and typing across all separations up to $15$\,m road baselines, decaying gracefully to $61\,\%$ at $25$\,m while holding unattributed perturbations below the trigger threshold.

\subsection{Gamified Accuracy and Pipeline Ablation Record}
The client gamification layer computes real-time confidence scores and structures participant rewards via an explicit XP progression: base crossing XP is augmented by a $2\times$ multi-band fusion bonus, $+15$\,XP for mesh-confirmed events, and a $1.25\times$ spatial consistency multiplier for nodes closest to the spatial consensus centroid. Because reward scheduling mirrors the operational likelihood distribution, incentive structures align with empirical measurement fidelity.

System architecture parameters were validated through complete-scene A/B testing rather than isolated row-level optimization:
\begin{itemize}
    \item \textbf{Numerical Optimizers:} Automated weight optimizers marginally boosted row-level $F_1$ scores ($0.173 \to 0.179$), yet degraded under full-scene simulation. An unconstrained optimizer converged to a degenerate state ($0\,\%$ human recall), while a constrained alternative reached $45.6\,\%$ scene accuracy, underperforming the shipped production heuristic ($46.7\,\%$).
    \item \textbf{Ablation of Rejected Strategies:} A/B testing exposed and eliminated several individually promising techniques: uniform-weight multi-band averaging (diluted BLE accuracy to $31.0\,\%$), noise-floor range extrapolation ($14$--$23$\,m phantom offsets), and naive sample-delta epoch comparison ($25$--$62\,\%$ false-positive rates).
\end{itemize}

\subsection{Limitations and untested conditions}
All results are model-based within a simulated test environment, meaning absolute metrics require empirical on-device re-validation; real-world accuracy is expected to degrade slightly due to richer physical multipath propagation. Within the modeled envelope, the system achieved identification up to approximately $12\text{ m}$ and detection to $\approx 16\text{ m}$, resolving into position bands rather than discrete points beyond $\approx 8\text{ m}$. Over this same $8\text{ m}$ span, epoch difference detection operated with zero false positives, while the mesh-extended configuration yielded a $20\text{--}40\text{ m}$ radius with zero false confirmations.

Several operational conditions remain untested and define future experiments on the same evaluation harness, pending on-device validation or per-environment refitting of $n$, $R_{1\text{m}}$, $A_0$, and $\sigma$. Indoor priorities include residential setups (2--5 links), office floors, corridor transitions, stairwell NLOS, RF-hostile settings, and high-density transit hubs with simultaneous crossings. Outdoor evaluations will cover urban street canyons, adverse weather (rain, wet ground, night), sparse park layouts ($30\text{--}50\text{ m}$ anchors), high-multipath parking lots, rural roads, and saturated events. Finally, cross-cutting investigations must address practical constraints such as device placement (hand, pocket, bag), hardware heterogeneity, battery-limited duty cycling, co-channel interference, and resilience against adversarial beacons and spoofed reports.

%---------------------------------------------------------------------------
% IV. CONCLUSIONS
%---------------------------------------------------------------------------
\section{Conclusions and Future Work}
Ambient Sense demonstrates that a commodity phone can perform multi-band passive
sensing end to end: fuse BLE and dual-band Wi-Fi features into entity typing,
convert measured dBm into a stated range with an explicit gate, detect change
between two instants of the field without inventing movement, count transit
across a two-node Bluetooth barrier, and accrue accuracy by agreeing with
neighbours --- while a gamification layer keeps voluntary contributors engaged.
Just as importantly, it reports where the physics withdraws its support: 6\,\%
typing accuracy at 68~m, 94\,\% rather than 100\,\% detection at 15~m, and a
noise floor below which ranging is not merely uncertain but refused.

\textbf{Use cases.} Because the sensing hardware is the handset, the marginal
cost of a new measurement point is approximately zero, and because the observer
carries no camera, the privacy argument is structural. The system supports
ambient-assisted living (presence and crossing counts without cameras or worn
devices), occupancy and flow analytics in workplaces and corridors, roadside
traffic studies (two phones become a typed barrier counter), and coverage
scouting, in which the RF field itself is the subject of the change-detection
epochs.

\textbf{Gamification and incentive extension.} The gamification layer will be
extended from a point-and-badge system to a quality-weighted, reputation-aware
reward engine: team meshes, seasonal ladders and accuracy-weighted seasons in
which the reliability bonus becomes a published, auditable score. The incentive
design is deliberately aligned with the measurement --- rewarding the devices
nearest the fused estimate is equivalent to rewarding honesty --- and draws on
established quality-aware and reputation-based mobile-crowdsensing
mechanisms~\cite{zhang2022,wang2018}.

\textbf{Summary.} The central contribution is a method of validation: every
constant that ships was produced by executing the harness and falsifying the
losers, so the paper can state, with equal precision, both the configuration
that shipped and the configuration the evidence threw away. Ambient Sense is a
working answer to a hard question --- what can a single, unmodified smartphone
held by an ordinary user actually deliver in the field --- and the answer
includes, as a first-class result, an explicit account of where it cannot.

%---------------------------------------------------------------------------
% REFERENCES
%---------------------------------------------------------------------------
\begin{thebibliography}{21}

\bibitem{wilson2010} J.~Wilson and N.~Patwari, ``Radio tomographic imaging with
wireless networks,'' \emph{IEEE Trans.\ Mobile Comput.}, vol.~9, no.~5,
pp.~621--632, May 2010, doi: 10.1109/TMC.2010.8.

\bibitem{januszkiewicz2018} {\L}.\,Januszkiewicz, ``Analysis of human body
shadowing effect on wireless sensor networks operating in the 2.4~GHz band,''
\emph{Sensors}, vol.~18, no.~10, art.~3412, Oct. 2018, doi:
10.3390/s18103412.

\bibitem{alippi2016} C.~Alippi, M.~Bocca, G.~Boracchi, N.~Patwari, and
M.~Roveri, ``RTI goes wild: Radio tomographic imaging for outdoor people
detection and localization,'' \emph{IEEE Trans.\ Mobile Comput.}, vol.~15,
no.~10, pp.~2585--2598, Oct. 2016, doi: 10.1109/TMC.2015.2504964.

\bibitem{li2024} Z.~Li, A.~Dubey, S.~Shen, N.~K.~Kundu, J.~Rao, and R.~Murch,
``Radio tomographic imaging with reconfigurable intelligent surfaces for
integrated sensing and communication,'' \emph{IEEE Trans.\ Wireless Commun.},
vol.~23, no.~11, pp.~15784--15797, Nov. 2024, doi: 10.1109/TWC.2024.3433011.

\bibitem{wu2024} H.~Wu, C.~Cheng, T.~Peng, H.~Zhou, and T.~Chen, ``Combining
transformer with a latent variable model for radio tomography based robust
device-free localization,'' \emph{Comput.\ Commun.}, vol.~224, art.~108022,
Oct. 2024, doi: 10.1016/j.comcom.2024.108022.

\bibitem{cnnrti2025} ``Passive localization based on radio tomography images
with CNN model utilizing WiFi RSSI,'' \emph{Sci.\ Rep.}, vol.~15, art.~15698,
2025, doi: 10.1038/s41598-025-99694-2.

\bibitem{wong2024} W.~E.~Wong, A.~H.~Wong, W.~Q.~Peh, and C.~K.~Tan, ``A
high-dimensional, multi-transceiver channel state information dataset for
enhanced human activity recognition,'' \emph{Data Brief}, vol.~55,
art.~110673, Aug. 2024, doi: 10.1016/j.dib.2024.110673.

\bibitem{sanchez2020} D.~S\'anchez-Rodr\'iguez, M.~A.~Quintana-Su\'arez,
I.~Alonso-Gonz\'alez, C.~Ley-Bosch, and J.~J.~S\'anchez-Medina, ``Fusion of
channel state information and received signal strength for indoor localization
using a single access point,'' \emph{Remote Sens.}, vol.~12, no.~12,
art.~1995, Jun. 2020, doi: 10.3390/rs12121995.

\bibitem{ramirez2021} R.~Ramirez, C.-Y.~Huang, C.-A.~Liao, P.-T.~Lin, H.-W.~Lin,
and S.-H.~Liang, ``A practice of BLE RSSI measurement for indoor positioning,''
\emph{Sensors}, vol.~21, no.~15, art.~5181, Jul. 2021, doi: 10.3390/s21155181.

\bibitem{milano2024} F.~Milano, H.~da Rocha, M.~Laracca, L.~Ferrigno,
A.~Esp\'iritu Santo, J.~Salvado, and V.~Paciello, ``BLE-based indoor
localization: Analysis of some solutions for performance improvement,''
\emph{Sensors}, vol.~24, no.~2, art.~376, Jan. 2024, doi: 10.3390/s24020376.

\bibitem{bai2020} L.~Bai, F.~Ciravegna, R.~Bond, and M.~Mulvenna, ``A low cost
indoor positioning system using Bluetooth Low Energy,'' \emph{IEEE Access},
vol.~8, pp.~136858--136871, 2020, doi: 10.1109/ACCESS.2020.3012342.

\bibitem{khan2023} D.~Khan and I.~W.-H.~Ho, ``CrossCount: Efficient device-free
crowd counting by leveraging transfer learning,'' \emph{IEEE Internet Things
J.}, vol.~10, no.~5, pp.~4049--4058, Mar. 2023.

\bibitem{zhou2020} R.~Zhou, Z.~Gong, X.~Lu, and Y.~Fu, ``WiFlowCount:
Device-free people flow counting by exploiting Doppler effect in commodity
WiFi,'' \emph{IEEE Syst.~J.}, vol.~14, no.~4, pp.~4919--4930, Dec. 2020.

\bibitem{dmip2023} ``Non-intrusive people counting and identification
simultaneously with commodity WiFi devices,'' in \emph{Proc.\ 6th Int.\ Conf.\
Digital Medicine and Image Processing (DMIP '23)}, 2023, pp.~62--68, doi:
10.1145/3637684.3637694.

\bibitem{zhang2022} J.~Zhang, X.~Li, Z.~Shi, \emph{et al.}, ``A reputation-based
and privacy-preserving incentive scheme for mobile crowd sensing: A deep
reinforcement learning approach,'' \emph{Wirel.\ Netw.}, 2023, doi:
10.1007/s11276-022-03111-9.

\bibitem{wang2018} J.~Wang, M.~Li, Y.~He, H.~Li, K.~Xiao, and C.~Wang, ``A
blockchain based privacy-preserving incentive mechanism in crowdsensing
applications,'' \emph{IEEE Access}, vol.~6, pp.~17545--17556, 2018, doi:
10.1109/ACCESS.2018.2815365.

\bibitem{casar2022} M.~C\"asar, T.~Pawelke, J.~Steffan, and G.~Terhorst, ``A
survey on Bluetooth Low Energy security and privacy,'' \emph{Comput.\ Netw.},
vol.~202, art.~108712, Jan. 2022, doi: 10.1016/j.comnet.2021.108712.

\bibitem{antonioli2020} D.~Antonioli, N.~O.~Tippenhauer, and K.~B.~Rasmussen,
``Key negotiation downgrade attacks on Bluetooth and Bluetooth Low Energy,''
\emph{ACM Trans.\ Priv.\ Secur.}, vol.~23, no.~3, art.~14, Jul. 2020, doi:
10.1145/3394497.

\bibitem{cai2024} H.~Cai, ``Securing billion Bluetooth devices leveraging
learning-based techniques,'' in \emph{Proc.\ AAAI Conf.\ Artif.\ Intell.\
(AAAI)}, vol.~38, no.~21, pp.~23731--23732, 2024, doi:
10.1609/aaai.v38i21.30544.

\bibitem{mobisys2024} ``Poster: Hybrid detection mechanism for spoofing
attacks in Bluetooth Low Energy networks,'' in \emph{Proc.\ 22nd Annu.\ Int.\
Conf.\ Mobile Syst., Appl., Services (MobiSys '24)}, 2024, pp.~710--711, doi:
10.1145/3643832.3661434.

\bibitem{hensler2019} C.~Hensler and P.~Tague, ``DEMO: Using Bluetooth Low Energy
spoofing to dispute device details,'' in \emph{Proc.\ ACM Conf.\ Security
Privacy Wireless Mobile Netw.\ (WiSec '19)}, Miami, FL, USA, 2019, pp.~341--343,
doi: 10.1145/3317549.3326321.

\end{thebibliography}

\end{document}
